\chapter{場景模擬}
\renewcommand{\baselinestretch}{10.0} %設定行距
\pagenumbering{arabic} %設定頁號阿拉伯數字
\setcounter{page}{7}  %設定頁數
\fontsize{14pt}{2.5pt}\sectionef
\section{摘要}
  完成球員及程式碼設定後，接著建立模擬場景，添加場地、球員、計時器、 LED 記分板和機械式轉盤記分板\\
\section{統整場景}
  在球場周圍設置感測器球碰到會返回中心,放入 8 名球員,以兩色分為兩隊。\\
\begin{figure}[hbt!]
\begin{center}
\includegraphics[height=8cm]{總2}
\caption{\Large 球場建立}\label{總2}
\end{center}
\end{figure}

  和計時器、LED計分板和機械式記分板。\\
\begin{figure}[hbt!]
\begin{center}
\includegraphics[height=15cm]{總}
\caption{\Large 球場全貌}\label{總}
\end{center}
\end{figure}
\section{ CoppeliaSim}
  CoppeliaSim，曾被稱為V-REP（Virtual Robot Experimentation Platform），是一個功能強大且多用途的機器人仿真軟體。它允許使用者在虛擬環境中創建、模擬和分析機器人系統。CoppeliaSim提供各種功能和功能，適用於各種機器人應用、研究和教育目的。\\
\begin{figure}[hbt!]
\begin{center}
\includegraphics[height=8cm]{CoppeliaSim}
\caption{\Large CoppeliaSim圖片}\label{CoppeliaSim}
\end{center}
\end{figure}
\renewcommand{\baselinestretch}{1} %設定行距